%% ============================================================
%% 论文修改稿：三大核心问题的完整修改
%% 问题3: 约束(8)(9)的非线性项线性化
%% 问题4: 约束(10)岸电容量约束的逻辑修正
%% 问题2: 补全列生成算法的完整技术细节
%% ============================================================
%% 使用说明：
%%   - 第一部分：替换原文 2.4.4 节中的约束 (4)-(6)
%%   - 第二部分：作为新的第 3 节插入（原第 3 节数值实验改为第 4 节）
%% ============================================================

\documentclass[11pt,a4paper]{article}
\usepackage[UTF8]{ctex}
\usepackage{amsmath,amssymb,amsthm}
\usepackage{algorithm,algorithmic}
\usepackage{booktabs}
\usepackage{geometry}
\geometry{left=2.5cm,right=2.5cm,top=2.5cm,bottom=2.5cm}

\newtheorem{proposition}{命题}
\newtheorem{remark}{注}
\newtheorem{definition}{定义}

\begin{document}

%% ============================================================
%% 第一部分：修正后的 MILP 约束条件
%% 替换原文 2.4.4 节中的约束 (4)-(9) 及相应文字
%% ============================================================

\section*{【修改一】2.4.4 节约束条件的修正版本}

\noindent\textbf{修改说明：}原约束 (8) 和 (9) 中存在二元变量乘积项 $w_{i\tau} \cdot x_{ik}$
和 $w_{i\tau} \cdot y_i$，使模型实质为非线性整数规划（MINLP），与文中声称的 MILP 不符。
原约束 (10) 对所有岸电桩聚合求和，未正确建立单桩独占约束。以下给出完整修正方案。

\bigskip
\hrule
\bigskip

\subsection*{2.4.3\quad 变量补充定义}

在原有变量基础上，为实现线性化，引入以下辅助变量：

\begin{table}[h]
\centering
\caption{补充决策变量定义}
\begin{tabular}{ll}
\toprule
变量 & 含义 \\
\midrule
$s_{ikt} \in \{0,1\}$ & 若船舶 $i$ 选择岸电桩 $k$ 且在时间段 $t$ 开始充电，则取 1 \\
$b_{it} \in \{0,1\}$ & 若船舶 $i$ 选择换电服务且在时间段 $t$ 开始换电，则取 1 \\
\bottomrule
\end{tabular}
\end{table}

\noindent\textbf{变量语义说明：} $s_{ikt}$ 联合编码了"船舶 $i$ 分配到岸电桩 $k$"
与"在时间段 $t$ 开始服务"两个决策，等价于原模型中 $x_{ik} \cdot w_{it}$ 的线性化替代。
类似地，$b_{it}$ 编码了"船舶 $i$ 选择换电"与"在时间段 $t$ 开始服务"，
等价于 $y_i \cdot w_{it}$。

\subsection*{2.4.4\quad 约束条件（修正版）}

\medskip
\noindent\textbf{(1) 服务模式选择约束}\quad 每艘船舶恰好选择一种能源服务模式：
%
\begin{equation}
  \sum_{k \in \mathcal{K}^{SP}} Q_{ik}\, x_{ik} \;+\; y_i \;+\; z_i \;=\; 1,
  \qquad \forall\, i \in \mathcal{N}.
  \tag{5}
\end{equation}

\medskip
\noindent\textbf{(2) 能源服务开始时刻约束}\quad
若船舶 $i$ 选择岸电或换电服务，则必须恰好在某一时间段启动：
%
\begin{equation}
  \sum_{t \in \mathcal{T}} w_{it}
  \;=\; \sum_{k \in \mathcal{K}^{SP}} x_{ik} \;+\; y_i,
  \qquad \forall\, i \in \mathcal{N}.
  \tag{6}
\end{equation}

\medskip
\noindent\textbf{(3) 服务持续时间定义}
%
\begin{equation}
  d_i^{SP} = \left\lceil \frac{E_i}{P^{SP}} \right\rceil, \qquad
  d_i^{BS} = T^{BS}.
  \tag{7}
\end{equation}

%% ===== 核心修正：约束 (8)(9) 的线性化 =====

\medskip
\noindent\textbf{(4) 岸电服务启动链接约束（线性化）}\quad
辅助变量 $s_{ikt}$ 联合表示"船舶 $i$ 分配至岸电桩 $k$ 并在时间段 $t$ 开始服务"。
以下约束组确保 $s_{ikt} = x_{ik} \cdot w_{it}$ 在二元域上成立，且为纯线性形式：
%
\begin{align}
  s_{ikt} &\;\leq\; x_{ik},
    & \forall\, i \in \mathcal{N},\; k \in \mathcal{K}^{SP},\; t \in \mathcal{T},
    \tag{8a} \\[4pt]
  s_{ikt} &\;\leq\; w_{it},
    & \forall\, i \in \mathcal{N},\; k \in \mathcal{K}^{SP},\; t \in \mathcal{T},
    \tag{8b} \\[4pt]
  s_{ikt} &\;\geq\; x_{ik} + w_{it} - 1,
    & \forall\, i \in \mathcal{N},\; k \in \mathcal{K}^{SP},\; t \in \mathcal{T},
    \tag{8c} \\[4pt]
  \sum_{t \in \mathcal{T}} s_{ikt} &\;=\; x_{ik},
    & \forall\, i \in \mathcal{N},\; k \in \mathcal{K}^{SP}.
    \tag{8d}
\end{align}

\noindent\textbf{约束解释：}
约束 (8a)--(8c) 是标准的 McCormick envelope（也称 Fortet 线性化），
用于精确线性化两个二元变量 $x_{ik}$ 与 $w_{it}$ 的乘积。
具体地：(8a) 和 (8b) 要求 $s_{ikt}$ 不超过任一因子；
(8c) 要求当两个因子同时为 1 时 $s_{ikt}$ 至少为 1。
对于二元变量，该松弛是紧的（tight），即 $s_{ikt} = x_{ik} \cdot w_{it}$ 恰好成立。
约束 (8d) 是冗余强化约束（valid inequality）：
若船舶 $i$ 被分配到岸电桩 $k$（$x_{ik}=1$），
则恰好存在一个时间段 $t$ 使得 $s_{ikt}=1$；
若 $x_{ik}=0$，则所有 $s_{ikt}=0$。
该约束虽然可由 (8a)--(8c) 与 (6) 联合推出，
但显式加入可显著增强 LP 松弛的紧致性。

\medskip
\noindent\textbf{(4$'$) 岸电占用变量定义}\quad
基于线性化变量 $s_{ikt}$，岸电桩 $k$ 在时间段 $t$ 被船舶 $i$ 占用的指示变量为：
%
\begin{equation}
  u_{ikt} \;=\; \sum_{\tau = \max\{0,\, t - d_i^{SP} + 1\}}^{t} s_{ik\tau},
  \qquad \forall\, i \in \mathcal{N},\; k \in \mathcal{K}^{SP},\; t \in \mathcal{T}.
  \tag{9}
\end{equation}

\noindent\textbf{说明：}与原约束 (8) 相比，此处将非线性项 $w_{i\tau} \cdot x_{ik}$
替换为已线性化的 $s_{ik\tau}$，从而使约束 (9) 成为纯线性定义式。
直观地，$u_{ikt} = 1$ 当且仅当船舶 $i$ 在某个时间段
$\tau \in [t - d_i^{SP} + 1,\; t]$ 开始在岸电桩 $k$ 上充电，
且充电过程尚未结束（即时间段 $t$ 仍处于服务窗口内）。

\medskip
\noindent\textbf{(5) 换电服务启动链接约束（线性化）}\quad
辅助变量 $b_{it}$ 联合表示"船舶 $i$ 选择换电并在时间段 $t$ 开始服务"，
约束组确保 $b_{it} = y_i \cdot w_{it}$：
%
\begin{align}
  b_{it} &\;\leq\; y_i,
    & \forall\, i \in \mathcal{N},\; t \in \mathcal{T},
    \tag{10a} \\[4pt]
  b_{it} &\;\leq\; w_{it},
    & \forall\, i \in \mathcal{N},\; t \in \mathcal{T},
    \tag{10b} \\[4pt]
  b_{it} &\;\geq\; y_i + w_{it} - 1,
    & \forall\, i \in \mathcal{N},\; t \in \mathcal{T},
    \tag{10c} \\[4pt]
  \sum_{t \in \mathcal{T}} b_{it} &\;=\; y_i,
    & \forall\, i \in \mathcal{N}.
    \tag{10d}
\end{align}

\noindent\textbf{(5$'$) 换电占用变量定义}
%
\begin{equation}
  v_{it} \;=\; \sum_{\tau = \max\{0,\, t - d_i^{BS} + 1\}}^{t} b_{i\tau},
  \qquad \forall\, i \in \mathcal{N},\; t \in \mathcal{T}.
  \tag{11}
\end{equation}

%% ===== 核心修正：约束 (10) 的逻辑修正 =====

\medskip
\noindent\textbf{(6) 资源容量约束（修正版）}\quad
岸电桩为独占性资源：每个岸电桩 $k$ 在任意时间段 $t$
至多被一艘船舶占用。换电设备为同质资源，
总占用量不超过设备数量 $K^{BS}$：
%
\begin{align}
  \sum_{i \in \mathcal{N}} u_{ikt} &\;\leq\; 1,
    & \forall\, k \in \mathcal{K}^{SP},\; t \in \mathcal{T},
    \tag{12} \\[4pt]
  \sum_{i \in \mathcal{N}} v_{it} &\;\leq\; K^{BS},
    & \forall\, t \in \mathcal{T}.
    \tag{13}
\end{align}

\noindent\textbf{修正说明：}原约束 (10) 写为
$\sum_{i \in \mathcal{N}} \sum_{k \in \mathcal{K}^{SP}} u_{ikt} \leq |\mathcal{K}^{SP}|$，
对所有岸电桩聚合求和。该形式仅约束了系统总使用量不超过桩的数量，
但未排除两艘船舶同时使用同一根桩的情况。
例如，若 $|\mathcal{K}^{SP}| = 2$，原约束允许桩 1 同时服务 2 艘船舶
而桩 2 空闲，因为 $2 + 0 = 2 \leq 2$，但这在物理上是不可行的。
修正后的约束 (12) 对每根岸电桩单独施加独占约束，正确刻画了"一桩一船"的物理限制。

\medskip
\noindent\textbf{(7) SIMOPS 完工时间约束}\quad
船舶离港时间由装卸作业和能源服务中较长者主导：
%
\begin{align}
  L_i &\;\geq\; A_i + T_{\text{cargo}}^{(i)},
    & \forall\, i \in \mathcal{N},
    \tag{14} \\[4pt]
  L_i &\;\geq\; \sum_{t \in \mathcal{T}} t \cdot w_{it}
    \;+\; d_i^{SP} \sum_{k \in \mathcal{K}^{SP}} x_{ik}
    \;+\; d_i^{BS}\, y_i,
    & \forall\, i \in \mathcal{N}.
    \tag{15}
\end{align}

\medskip
\noindent\textbf{(8) 延误量线性化}
%
\begin{align}
  \delta_i &\;\geq\; L_i - D_i,
    & \forall\, i \in \mathcal{N},
    \tag{16} \\[4pt]
  \delta_i &\;\geq\; 0,
    & \forall\, i \in \mathcal{N}.
    \tag{17}
\end{align}

\medskip
\noindent\textbf{(9) 变量域约束}
%
\begin{equation}
  x_{ik},\; y_i,\; z_i,\; w_{it},\; s_{ikt},\; b_{it},\;
  u_{ikt},\; v_{it} \in \{0,1\},
  \qquad L_i,\; \delta_i \geq 0.
  \tag{18}
\end{equation}

\begin{remark}[模型规模分析]
线性化引入的额外变量 $s_{ikt}$ 和 $b_{it}$ 使得变量总数增加
$O(|\mathcal{N}| \cdot |\mathcal{K}^{SP}| \cdot |\mathcal{T}|)$ 和 $O(|\mathcal{N}| \cdot |\mathcal{T}|)$。
然而，由于 McCormick 线性化对二元变量是紧的，
修正后模型的 LP 松弛界不弱于原非线性模型的连续松弛，
从而不会降低求解效率。实际上，显式的线性化约束
反而有助于商业求解器（如 Gurobi、CPLEX）利用预处理和割平面技术，
往往比在求解过程中隐式处理双线性项更加高效。
\end{remark}


\newpage
%% ============================================================
%% 第二部分：列生成算法的完整技术细节
%% 作为新的第 3 节插入
%% ============================================================

\section*{【修改二】新增第 3 节：求解方法}

\setcounter{section}{2}
\setcounter{equation}{18}

\section{求解方法：基于列生成的分解框架}

第 2 节建立的混合整数线性规划模型（以下简称紧凑模型，Compact MILP）
在理论上可由商业求解器直接求解。
然而，时间展开变量 $w_{it}$、$s_{ikt}$、$u_{ikt}$ 等的数量随船舶规模 $|\mathcal{N}|$、
岸电桩数 $|\mathcal{K}^{SP}|$ 及时间段数 $|\mathcal{T}|$ 呈多项式增长，
导致大规模实例下的直接求解面临严峻的计算瓶颈。
例如，在 $|\mathcal{N}| = 100$、$|\mathcal{K}^{SP}| = 2$、$|\mathcal{T}| = 32$
的典型配置下，仅 $s_{ikt}$ 变量即达 $100 \times 2 \times 32 = 6{,}400$ 个，
加上其他时间展开变量，总二元变量数可超过 $20{,}000$。

为此，本节提出一种基于列生成（Column Generation, CG）的分解求解框架。
其核心思想是：将每艘船舶的所有可行服务方案（包含服务模式选择、
服务开始时刻及相应资源占用轮廓）预先定义为"列"（column），
通过限制性主问题（Restricted Master Problem, RMP）
与定价子问题（Pricing Subproblem, PP）的交替求解，
迭代生成并选择高质量的服务方案，
避免显式枚举全部时间展开变量。

\subsection{列的定义与可行服务方案}

\begin{definition}[服务方案]
对于船舶 $i \in \mathcal{N}$，一个\textbf{可行服务方案} $p \in \mathcal{P}_i$
由以下元素构成：
\begin{itemize}
  \item 服务模式 $m_p \in \{SP, BS, AE\}$；
  \item 若 $m_p = SP$：所分配的岸电桩 $k_p \in \mathcal{K}^{SP}$
        且满足兼容性 $Q_{ik_p} = 1$，以及服务开始时间段 $t_p \in \mathcal{T}$
        且 $t_p \geq A_i$；
  \item 若 $m_p = BS$：服务开始时间段 $t_p \in \mathcal{T}$
        且 $t_p \geq A_i$；
  \item 若 $m_p = AE$：无需额外参数（无资源占用）。
\end{itemize}
\end{definition}

\noindent
对于每个服务方案 $p \in \mathcal{P}_i$，其参数可由如下方式计算：

\medskip
\noindent\textbf{直接成本：}
\begin{equation}
  c_{ip} =
  \begin{cases}
    C^{SP} \cdot E_i, & \text{若 } m_p = SP, \\
    C^{BS}, & \text{若 } m_p = BS, \\
    C^{AE} \cdot E_i, & \text{若 } m_p = AE.
  \end{cases}
  \tag{19}
\end{equation}

\noindent\textbf{完工时间与延误成本：}
\begin{align}
  L_{ip} &= \max\!\Big\{A_i + T_{\text{cargo}}^{(i)},\;
    t_p + d_{ip}\Big\},
  \tag{20} \\[4pt]
  \delta_{ip} &= \max\!\big\{L_{ip} - D_i,\; 0\big\},
  \tag{21} \\[4pt]
  \bar{c}_{ip} &= c_{ip} + C_{\text{delay}}^{(i)} \cdot \delta_{ip},
  \tag{22}
\end{align}
%
其中 $d_{ip}$ 为方案 $p$ 的服务持续时间：
若 $m_p = SP$ 则 $d_{ip} = d_i^{SP}$；
若 $m_p = BS$ 则 $d_{ip} = d_i^{BS}$；
若 $m_p = AE$ 则 $d_{ip} = 0$。
$\bar{c}_{ip}$ 为方案 $p$ 的\textbf{广义成本}，
涵盖能源直接成本与延误惩罚。

\medskip
\noindent\textbf{资源占用轮廓：}
对于方案 $p \in \mathcal{P}_i$，定义其在时间段 $t$ 对岸电桩 $k$ 和换电设备的占用：
%
\begin{align}
  \alpha_{ipkt} &=
  \begin{cases}
    1, & \text{若 } m_p = SP,\; k = k_p,\;
         t_p \leq t \leq t_p + d_i^{SP} - 1, \\
    0, & \text{否则},
  \end{cases}
  \tag{23} \\[6pt]
  \beta_{ipt} &=
  \begin{cases}
    1, & \text{若 } m_p = BS,\;
         t_p \leq t \leq t_p + d_i^{BS} - 1, \\
    0, & \text{否则}.
  \end{cases}
  \tag{24}
\end{align}

\subsection{限制性主问题（RMP）}

设 $\hat{\mathcal{P}}_i \subseteq \mathcal{P}_i$ 为当前迭代中船舶 $i$
可用的服务方案子集。
引入决策变量 $\lambda_{ip} \in [0,1]$（LP 松弛域），
表示船舶 $i$ 是否采用方案 $p$。
限制性主问题的 LP 松弛形式为：

\begin{align}
  [\text{RMP}] \qquad
  \min \quad & \sum_{i \in \mathcal{N}} \sum_{p \in \hat{\mathcal{P}}_i}
    \bar{c}_{ip}\, \lambda_{ip}
  \tag{25} \\[6pt]
  \text{s.t.} \quad
  & \sum_{p \in \hat{\mathcal{P}}_i} \lambda_{ip} = 1,
    \qquad \forall\, i \in \mathcal{N},
    \qquad [\pi_i]
  \tag{26} \\[4pt]
  & \sum_{i \in \mathcal{N}} \sum_{p \in \hat{\mathcal{P}}_i}
    \alpha_{ipkt}\, \lambda_{ip} \leq 1,
    \qquad \forall\, k \in \mathcal{K}^{SP},\; t \in \mathcal{T},
    \qquad [\mu_{kt}]
  \tag{27} \\[4pt]
  & \sum_{i \in \mathcal{N}} \sum_{p \in \hat{\mathcal{P}}_i}
    \beta_{ipt}\, \lambda_{ip} \leq K^{BS},
    \qquad \forall\, t \in \mathcal{T},
    \qquad [\nu_t]
  \tag{28} \\[4pt]
  & 0 \leq \lambda_{ip} \leq 1,
    \qquad \forall\, i \in \mathcal{N},\; p \in \hat{\mathcal{P}}_i.
  \tag{29}
\end{align}

\noindent\textbf{约束说明：}
\begin{itemize}
  \item 约束 (26)（凸性约束）：每艘船舶恰好选择一个服务方案。
        对偶变量 $\pi_i$ 为自由变量，表示船舶 $i$ 的"影子价格"。
  \item 约束 (27)（岸电独占约束）：每根岸电桩在每个时间段最多服务一艘船。
        对偶变量 $\mu_{kt} \geq 0$，反映时段 $t$ 上岸电桩 $k$ 的拥挤价格。
  \item 约束 (28)（换电容量约束）：换电设备总占用不超过容量上限。
        对偶变量 $\nu_t \geq 0$，反映换电资源在时段 $t$ 的稀缺程度。
\end{itemize}

\begin{remark}[与紧凑模型的等价性]
当 $\hat{\mathcal{P}}_i = \mathcal{P}_i$（即包含所有可行方案）时，
RMP 的整数最优解与紧凑模型 (4)--(18) 的最优解等价。
这是因为每个整数可行解
$(\mathbf{x}, \mathbf{y}, \mathbf{z}, \mathbf{w})$
可唯一映射为某个方案选择 $\lambda_{ip} = 1$，反之亦然。
列生成的核心在于：无需显式枚举 $\mathcal{P}_i$，
而是通过定价子问题按需生成有价值的列。
\end{remark}

\subsection{定价子问题（Pricing Subproblem）}

在每次 RMP 求解后，利用对偶最优解 $(\boldsymbol{\pi}^*,\, \boldsymbol{\mu}^*,\, \boldsymbol{\nu}^*)$
检验是否存在尚未加入 RMP 的列具有负的缩减成本（reduced cost）。
对于船舶 $i$，方案 $p$ 的缩减成本定义为：
%
\begin{equation}
  \bar{\sigma}_{ip}
  \;=\; \bar{c}_{ip}
  \;-\; \pi_i^*
  \;-\; \sum_{k \in \mathcal{K}^{SP}} \sum_{t \in \mathcal{T}}
    \mu_{kt}^* \cdot \alpha_{ipkt}
  \;-\; \sum_{t \in \mathcal{T}} \nu_t^* \cdot \beta_{ipt}.
  \tag{30}
\end{equation}

\noindent
若 $\bar{\sigma}_{ip} < 0$，则将方案 $p$ 加入 $\hat{\mathcal{P}}_i$
可改善 RMP 的目标值。

对于每艘船舶 $i$，定价子问题为：
\begin{equation}
  [\text{PP}_i] \qquad
  \bar{\sigma}_i^* \;=\;
  \min_{p \in \mathcal{P}_i} \; \bar{\sigma}_{ip}.
  \tag{31}
\end{equation}

\noindent
将 (30) 中各项展开，$\text{PP}_i$ 可按服务模式分解为三个独立的子问题：

\medskip
\noindent\textbf{(a) 岸电定价子问题}
($m_p = SP$)：对于每对 $(k, t)$，其中 $k \in \mathcal{K}^{SP}$
且 $Q_{ik} = 1$，$t \in [A_i,\; |\mathcal{T}| - d_i^{SP} + 1]$，
计算：
%
\begin{equation}
  \sigma_i^{SP}(k, t)
  \;=\; \underbrace{C^{SP} \cdot E_i + C_{\text{delay}}^{(i)} \cdot \delta_i(t, d_i^{SP})}_{
    \text{广义成本 } \bar{c}_{ip}}
  \;-\; \pi_i^*
  \;-\; \underbrace{\sum_{\tau = t}^{t + d_i^{SP} - 1} \mu_{k\tau}^*}_{
    \text{岸电资源对偶价格}},
  \tag{32}
\end{equation}
%
其中 $\delta_i(t, d) = \max\!\big\{\max\{t + d,\; A_i + T_{\text{cargo}}^{(i)}\} - D_i,\; 0\big\}$
为在时间段 $t$ 开始、持续 $d$ 时段后的延误量。
最优岸电方案为：
\begin{equation}
  (k_i^*, t_i^{SP*}) = \arg\min_{k, t} \; \sigma_i^{SP}(k, t).
  \tag{33}
\end{equation}

\noindent
该子问题的计算复杂度为 $O(|\mathcal{K}^{SP}| \cdot |\mathcal{T}| \cdot d_{\max}^{SP})$，
其中 $d_{\max}^{SP} = \max_{i} d_i^{SP}$。
利用对偶价格的前缀和预计算
$M_{k}(t) = \sum_{\tau=0}^{t} \mu_{k\tau}^*$，
可将每对 $(k, t)$ 的岸电对偶价格求和简化为
$M_{k}(t + d_i^{SP} - 1) - M_{k}(t - 1)$，
使复杂度降至 $O(|\mathcal{K}^{SP}| \cdot |\mathcal{T}|)$。

\medskip
\noindent\textbf{(b) 换电定价子问题}
($m_p = BS$)：对于每个 $t \in [A_i,\; |\mathcal{T}| - d_i^{BS} + 1]$，
计算：
%
\begin{equation}
  \sigma_i^{BS}(t)
  \;=\; C^{BS} + C_{\text{delay}}^{(i)} \cdot \delta_i(t, d_i^{BS})
  \;-\; \pi_i^*
  \;-\; \sum_{\tau = t}^{t + d_i^{BS} - 1} \nu_{\tau}^*.
  \tag{34}
\end{equation}
%
类似地利用前缀和，复杂度为 $O(|\mathcal{T}|)$。

\medskip
\noindent\textbf{(c) 褐色服务定价}
($m_p = AE$)：无资源占用，缩减成本为：
%
\begin{equation}
  \sigma_i^{AE}
  \;=\; C^{AE} \cdot E_i
  \;+\; C_{\text{delay}}^{(i)} \cdot \max\!\big\{A_i + T_{\text{cargo}}^{(i)} - D_i,\; 0\big\}
  \;-\; \pi_i^*.
  \tag{35}
\end{equation}
%
这是一个常数（不依赖时间维度），复杂度为 $O(1)$。

\medskip
\noindent\textbf{汇总：}船舶 $i$ 的最优缩减成本为：
\begin{equation}
  \bar{\sigma}_i^* = \min\!\Big\{
    \min_{k,t} \sigma_i^{SP}(k,t),\;\;
    \min_{t} \sigma_i^{BS}(t),\;\;
    \sigma_i^{AE}
  \Big\}.
  \tag{36}
\end{equation}

\noindent
定价子问题的\textbf{总复杂度}为
$O\!\big(|\mathcal{N}| \cdot (|\mathcal{K}^{SP}| + 1) \cdot |\mathcal{T}|\big)$，
在典型参数下（$|\mathcal{N}| = 500$, $|\mathcal{K}^{SP}| = 2$,
$|\mathcal{T}| = 32$）仅需约 $48{,}000$ 次基本运算，可在毫秒级完成。

\subsection{初始列池构造}

列生成算法需要一个可行的初始列池以启动 RMP 求解。
本文采用以下策略构造初始列池 $\hat{\mathcal{P}}_i^{(0)}$：

\begin{enumerate}
  \item \textbf{褐色保底列：}对每艘船舶 $i$，加入褐色服务方案
        $p_i^{AE}$（$m_p = AE$）。由于褐色服务不消耗任何港口资源，
        该方案始终可行，从而保证 RMP 在任意迭代中均有可行解。
        该设计与运输调度中"外包兜底"的建模思想一致~\cite{Cordeau2015}。
  \item \textbf{贪心岸电列：}对每艘兼容岸电的船舶 $i$
        （$\exists\, k: Q_{ik} = 1$），按到达时间排序，
        贪心分配最早可用的岸电桩和最早可行开始时间。
  \item \textbf{即时换电列：}对每艘船舶 $i$，
        加入在到达时刻立即换电的方案（$t_p = A_i$，$m_p = BS$）。
\end{enumerate}

\noindent
该初始化策略保证了两个性质：
(i) RMP 在首次迭代即可行（由褐色列保证）；
(ii) 初始解质量不过差（贪心列和即时列提供了合理的起始点），
有助于加速后续列生成的收敛。

\subsection{列生成迭代流程与终止准则}

算法~\ref{alg:cg} 给出了列生成框架的完整流程。

\begin{algorithm}[htbp]
\caption{列生成求解框架}
\label{alg:cg}
\begin{algorithmic}[1]
\REQUIRE 问题实例 $(\mathcal{N}, \mathcal{K}^{SP}, \mathcal{K}^{BS}, \mathcal{T},
  \text{参数})$；收敛阈值 $\epsilon > 0$；最大迭代数 $I_{\max}$
\ENSURE 近最优服务方案分配
\STATE 构造初始列池 $\hat{\mathcal{P}}_i^{(0)}$，$\forall\, i \in \mathcal{N}$
  \hfill \textit{// 第 3.4 节}
\STATE $\ell \leftarrow 0$
\REPEAT
  \STATE $\ell \leftarrow \ell + 1$
  \STATE 求解 RMP 的 LP 松弛，获得原始解 $\boldsymbol{\lambda}^*$
    和对偶解 $(\boldsymbol{\pi}^*,\, \boldsymbol{\mu}^*,\, \boldsymbol{\nu}^*)$
    \hfill \textit{// 第 3.2 节}
  \STATE $\textit{improved} \leftarrow \textsc{False}$
  \FOR{each $i \in \mathcal{N}$}
    \STATE 求解 $\text{PP}_i$：计算 $\bar{\sigma}_i^*$
      及对应最优方案 $p_i^*$
      \hfill \textit{// 第 3.3 节}
    \IF{$\bar{\sigma}_i^* < -\epsilon$}
      \STATE $\hat{\mathcal{P}}_i \leftarrow \hat{\mathcal{P}}_i \cup \{p_i^*\}$
      \STATE $\textit{improved} \leftarrow \textsc{True}$
    \ENDIF
  \ENDFOR
\UNTIL{$\textit{improved} = \textsc{False}$ \textbf{or} $\ell \geq I_{\max}$}
\STATE 在最终列池 $\hat{\mathcal{P}}_i$ 上
  将 $\lambda_{ip}$ 恢复为整数变量，求解整数 RMP
\RETURN 最优整数解对应的服务方案分配
\end{algorithmic}
\end{algorithm}

\subsection{收敛性与最优性分析}

\begin{proposition}[LP 松弛最优性]
若算法~\ref{alg:cg} 因 $\textit{improved} = \textsc{False}$ 终止
（即所有船舶的最小缩减成本 $\bar{\sigma}_i^* \geq -\epsilon$），
则 RMP 的当前 LP 最优值 $Z_{LP}^*$ 满足：
\begin{equation}
  Z_{LP}^* \;\leq\; Z_{LP}^{\text{full}} + |\mathcal{N}| \cdot \epsilon,
  \tag{37}
\end{equation}
其中 $Z_{LP}^{\text{full}}$ 为完整列集 $\mathcal{P}_i$ 上 LP 松弛的最优值。
当 $\epsilon \to 0$ 时，$Z_{LP}^* = Z_{LP}^{\text{full}}$，
即列生成精确求解了 LP 松弛。
\end{proposition}

\begin{proof}
这是列生成框架的标准收敛性结果。
当所有定价子问题的最优缩减成本非负时，
由线性规划的对偶理论，当前 RMP 对偶解对于完整列集也是可行的，
其对偶目标值即为 LP 松弛的最优值。
由强对偶定理，原始最优值 $Z_{LP}^*$ 等于对偶最优值，
从而等于完整列集上的 LP 最优值。
详见 Desaulniers et al.~\cite{Desaulniers2005}
第 2 章的标准证明。
\end{proof}

\begin{proposition}[整数最优性条件]
\label{prop:integer}
设 CG 收敛后 LP 最优值为 $Z_{LP}^*$，
最终整数 RMP 的最优值为 $Z_{IP}^*$。则：
\begin{equation}
  Z_{LP}^* \;\leq\; Z_{IP}^{\text{opt}} \;\leq\; Z_{IP}^*,
  \tag{38}
\end{equation}
其中 $Z_{IP}^{\text{opt}}$ 为紧凑模型 (4)--(18) 的全局整数最优值。
当 $Z_{LP}^* = Z_{IP}^*$ 时，
列生成获得了全局整数最优解。
\end{proposition}

\begin{proof}
第一个不等式成立是因为 LP 松弛提供整数规划的下界。
第二个不等式成立是因为限制列集上的整数最优解
是完整列集上的一个可行整数解。
当两个界重合时，最优性得证。
\end{proof}

\begin{remark}[实际求解中的整数恢复]
在实验中，当 $|\mathcal{N}| \leq 100$ 时，
CG 收敛后在最终列池上求解整数 RMP，
所得目标值与紧凑 MILP 的最优值完全一致
（$Z_{LP}^* = Z_{IP}^*$，见第 4 节表 7），
验证了命题~\ref{prop:integer} 中的充分条件在中等规模下普遍成立。
当 $|\mathcal{N}| > 100$ 时，
本文在每次定价迭代中仅保留缩减成本最优的 $K$ 列
（$K = 3$），以控制列池规模，
并设置最大迭代数 $I_{\max} = 20$ 和总时间上限 60 秒。
此时解可能不具有严格最优性保证，
但由 (38) 可计算最优性间隙上界
$\text{gap} = (Z_{IP}^* - Z_{LP}^*) / Z_{LP}^*$，
为解质量提供了可量化的评估依据。
\end{remark}

\subsection{算法复杂度总结}

表~\ref{tab:complexity} 汇总了列生成框架各组件的计算复杂度。

\begin{table}[htbp]
\centering
\caption{列生成框架计算复杂度分析}
\label{tab:complexity}
\begin{tabular}{lll}
\toprule
组件 & 单次复杂度 & 说明 \\
\midrule
RMP 求解 & $O(\text{LP}(|\hat{\mathcal{P}}|,\, m))$
  & $|\hat{\mathcal{P}}|$ 为列池大小，$m$ 为约束数 \\
岸电定价 & $O(|\mathcal{N}| \cdot |\mathcal{K}^{SP}| \cdot |\mathcal{T}|)$
  & 利用前缀和优化 \\
换电定价 & $O(|\mathcal{N}| \cdot |\mathcal{T}|)$
  & 利用前缀和优化 \\
褐色定价 & $O(|\mathcal{N}|)$
  & 常数时间 / 船 \\
整数恢复 & $O(\text{MIP}(|\hat{\mathcal{P}}|,\, m))$
  & 在收敛列池上求解 \\
\midrule
总迭代数 & $O(I_{\max})$
  & 典型值 $I_{\max} = 20$ \\
\bottomrule
\end{tabular}
\end{table}

\noindent
与直接求解紧凑 MILP 相比，列生成的核心优势在于：
(i) RMP 的规模由活跃列数而非全部时间展开变量决定，
大幅减少了 LP 求解的计算负担；
(ii) 定价子问题按船舶独立分解，可并行求解；
(iii) 通过 LP 对偶信息引导列的生成方向，
避免了对 $O(|\mathcal{N}| \cdot |\mathcal{K}^{SP}| \cdot |\mathcal{T}|)$
可行方案的显式枚举。

%% ============================================================
%% 参考文献（仅列出本修改新增引用的条目）
%% ============================================================

\begin{thebibliography}{99}

\bibitem{Desaulniers2005}
G.~Desaulniers, J.~Desrosiers, and M.M.~Solomon.
\textit{Column Generation}.
Springer, New York, 2005.

\bibitem{Cordeau2015}
J.-F.~Cordeau, G.~Laporte, M.W.P.~Savelsbergh, and D.~Vigo.
Transportation and logistics.
\textit{Handbooks in Operations Research and Management Science}, 14:1--43, 2015.

\bibitem{Barnhart1998}
C.~Barnhart, E.L.~Johnson, G.L.~Nemhauser, M.W.P.~Savelsbergh, and P.H.~Vance.
Branch-and-price: Column generation for solving huge integer programs.
\textit{Operations Research}, 46(3):316--329, 1998.

\bibitem{Lubbecke2005}
M.E.~L\"ubbecke and J.~Desrosiers.
Selected topics in column generation.
\textit{Operations Research}, 53(6):1007--1023, 2005.

\end{thebibliography}

\end{document}
